\documentclass{report}
\usepackage[utf8]{inputenc}
\usepackage[T1]{fontenc}
\usepackage[french]{babel}
\usepackage{geometry}
\usepackage{fancyhdr}
\setlength{\headheight}{13pt}
\usepackage{lastpage}
\usepackage{titlesec}
\usepackage{listings}
\usepackage{enumitem}
\usepackage{color}
\usepackage{xcolor}
\usepackage{soul}
\usepackage{graphicx}
\usepackage[utf8]{inputenc}
\usepackage{amsmath}
\usepackage{hyperref}
\hypersetup{hidelinks}
\usepackage{tikz}
\usetikzlibrary{positioning, shapes.geometric, arrows.meta, calc, fit, backgrounds, shadows}
\usepackage{rotating}
\usepackage{longtable}



\newcommand{\DocName}{OPTIMISATION DE MODELES NEURONAUX A OSCILLATEURS COUPLÉS - LIVRABLE 1}
\newcommand{\DocAuthors}{\bsc{Da Silva} Tristan - \bsc{Battaglia} Charles-Antoine \\ \bsc{Botté-Magalhaes} Jules - \bsc{Gauthier} Keanu - \bsc{Fraysse} Ianis}

\definecolor{dkgreen}{rgb}{0,0.6,0}
\definecolor{gray}{rgb}{0.5,0.5,0.5}
\definecolor{mauve}{rgb}{0.58,0,0.82}
\definecolor{addrred}{rgb}{0.93,0.43,0.43}
\definecolor{commentgray}{rgb}{0.45,0.45,0.45}

% Palette du schéma conceptuel
\definecolor{hubbg}{RGB}{23,48,77}
\definecolor{hubborder}{RGB}{23,48,77}
\definecolor{hubrule}{RGB}{150,178,206}
\definecolor{cardink}{RGB}{33,41,54}
\definecolor{cardline}{RGB}{160,172,186}
\definecolor{whyborder}{RGB}{31,92,139}
\definecolor{whybg}{RGB}{243,248,252}
\definecolor{whatborder}{RGB}{34,116,80}
\definecolor{whatbg}{RGB}{243,250,246}
\definecolor{howborder}{RGB}{181,84,39}
\definecolor{howbg}{RGB}{253,247,243}
\definecolor{outborder}{RGB}{104,72,140}
\definecolor{outbg}{RGB}{249,246,253}
\definecolor{defiborder}{RGB}{166,54,45}
\definecolor{defibg}{RGB}{253,245,244}
\definecolor{whoborder}{RGB}{71,85,105}
\definecolor{whobg}{RGB}{246,248,250}

% Outils de mise en forme des cartes du schéma conceptuel
\newlist{cbul}{itemize}{1}
\setlist[cbul]{leftmargin=3.4mm, labelsep=1.7mm, itemsep=2.6pt, topsep=0pt,
               parsep=0pt, label=\textcolor{cardaccent}{\rule[0.55ex]{1.4mm}{0.6pt}}}
\colorlet{cardaccent}{cardline}
\newcommand{\cardsetup}[1]{%
    \colorlet{cardaccent}{#1}\raggedright
    \hyphenpenalty=10000\exhyphenpenalty=10000\relax}
\newcommand{\cardtitle}[2]{%
    {\bfseries\color{#1}#2}\\[2.5pt]
    {\color{#1!35!white}\rule{\linewidth}{0.7pt}}\\[3.5pt]}

\lstdefinelanguage{MicroAsm}{
    sensitive=true,
    morecomment=[l]{\#},
    moredelim=[s][\color{addrred}]{@}{:}
}

\lstdefinelanguage{HexProg}{
    sensitive=true,
    morecomment=[l]{;}
}

\lstdefinelanguage{AsmStd}{
    sensitive=true,
    morecomment=[l]{;},
    morecomment=[l]{\#},
    morekeywords={
        MOV,ADD,SUB,MUL,DIV,INC,DEC,CMP,AND,OR,XOR,NOT,
        JMP,JE,JNE,JG,JL,JGE,JLE,CALL,RET,PUSH,POP,
        LDA,STA,LOAD,STORE,INPUT,OUTPUT,HALT,NOP
    }
}

\lstset{frame=tb,
    language=MicroAsm,
    aboveskip=3mm,
    belowskip=3mm,
    showstringspaces=false,
    columns=flexible,
    basicstyle={\small\ttfamily},
    numbers=none,
    numberstyle=\tiny\color{gray},
    keywordstyle=\color{blue},
    commentstyle=\color{commentgray}, 
    breaklines=true,
    breakatwhitespace=true,
    tabsize=3,
   }

\lstset{breaklines=true}

\geometry{a4paper, margin=1in}

% Suppression de l'affichage du mot "Chapitre"
\titleformat{\chapter}[display]
    {\normalfont\bfseries\LARGE}
    {}
    {0pt}
    {}

\definecolor{Light}{gray}{.90}
\sethlcolor{Light}

\let\OldTexttt\texttt
\renewcommand{\texttt}[1]{\OldTexttt{\hl{#1}}}

\renewcommand{\thechapter}{\arabic{chapter}}
\renewcommand{\thesection}{\thechapter.\arabic{section}}
\renewcommand{\thesubsection}{\thesection.\arabic{subsection}}

\fancyhf{}
\fancyhead[L]{\DocName}
\fancyfoot[L]{\DocAuthors}
\fancyfoot[R]{\thepage / \pageref{LastPage}}
\renewcommand{\headrulewidth}{0 pt}
\renewcommand{\footrulewidth}{0.2pt}
\pagestyle{fancy}
\fancypagestyle{plain}{
    \fancyhf{}
    \fancyfoot[R]{\thepage / \pageref{LastPage}}
    \fancyfoot[L]{\DocAuthors}
    \fancyhead[L]{\DocName}
}

\graphicspath{ {./img} }

\title{\DocName}
\author{\DocAuthors}
\date{\textbf{JUNIA - ISEN - AP4} \url{https://github.com/KeanuGauthier/Socle-Modele-Neuronaux-RD.git}}

\newcounter{questioncounter}
\setcounter{questioncounter}{0}

\begin{document}
\maketitle

\chapter*{Lexique}
\addcontentsline{toc}{chapter}{Lexique}

\begin{tabular}{@{}p{4cm}p{11cm}@{}}
    \textbf{Benchmark} & Test de performance comparant les temps d'exécution entre deux architectures, ici CPU et GPU. \\[6pt]
    \textbf{Canaux ioniques} & Pores membranaires permettant le passage sélectif d'ions (sodium, potassium) à travers la membrane du neurone. \\[6pt]
    \textbf{Connectome} & Carte complète des connexions synaptiques entre les neurones d'un réseau. \\[6pt]
    \textbf{CPU} & \textit{Central Processing Unit}, processeur généraliste traitant les instructions de façon séquentielle. \\[6pt]
    \textbf{EDO} & Équation Différentielle Ordinaire, équation décrivant l'évolution d'une variable par rapport au temps. \\[6pt]
    \textbf{GPU} & \textit{Graphics Processing Unit}, processeur massivement parallèle conçu pour effectuer un grand nombre de calculs simultanément. \\[6pt]
    \textbf{Hodgkin-Huxley} & Modèle biophysique décrivant l'évolution du potentiel membranaire d'un neurone à partir de canaux ioniques sodium et potassium. \\[6pt]
    \textbf{Potentiel membranaire} & Différence de potentiel électrique entre l'intérieur et l'extérieur de la membrane d'un neurone. \\[6pt]
    \textbf{RTX 5090} & Carte graphique haut de gamme de NVIDIA, équipée d'environ 21\,760 cœurs de calcul parallèle. \\[6pt]
    \textbf{Speedup} & Rapport entre le temps d'exécution sur CPU et sur GPU, mesurant le gain apporté par le GPU. \\[6pt]
    \textbf{Synapse} & Jonction entre deux neurones assurant la transmission du signal électrique d'un neurone à un autre. \\[6pt]
\end{tabular}

\tableofcontents

\chapter{Contexte et motivation}

Ce projet s'inscrit dans le cadre du laboratoire de recherche de M.~Antoine Pirog, encadrant du projet R\&D au semestre 7 à Junia. Le laboratoire travaille sur la simulation de réseaux de neurones biophysiques, c'est-à-dire des modèles qui reproduisent le comportement électrique de populations de neurones interconnectés.

Le problème auquel ce type de recherche fait face est simple : la simulation de réseaux de neurones réalistes est coûteuse en temps de calcul. Sur \textbf{CPU}, un réseau de quelques milliers de neurones peut prendre des heures à simuler. Cela limite la taille des réseaux qu'on peut étudier, et donc les questions scientifiques qu'on peut explorer.

Junia dispose de serveurs équipés de cartes graphiques \textbf{RTX~5090}. La question posée est donc : ces \textbf{GPU} peuvent-ils apporter la puissance de calcul nécessaire pour rendre la simulation de grands réseaux de neurones tractable au quotidien pour le laboratoire ?

La réponse n'est pas évidente, et elle n'est pas garantie. Le \textbf{GPU} est théoriquement beaucoup plus puissant, mais notre problème, un système d'\textbf{EDO} couplées avec des connexions irrégulières, n'est pas forcément un cas idéal pour le parallélisme massif. Il est tout à fait possible que la conclusion soit : non, sur ce type de simulation, le \textbf{GPU} n'apporte pas un gain suffisant pour justifier le changement d'architecture. Le but du projet est d'apporter une réponse documentée et mesurée à cette question, pas d'arriver à une conclusion prédéterminée.

C'est dans ce contexte que le sujet \og Optimisation de modèles neuronaux à oscillateurs couplés \fg{} nous a été confié, avec un socle de code déjà existant : modèle de \textbf{Hodgkin-Huxley}, \textbf{synapses}, solveur \textbf{CPU} de référence. On part de cette base pour évaluer l'intérêt du \textbf{GPU}.

\chapter{Présentation du sujet}

\section{Énoncé et objectif}

Ce projet a pour objectif d'évaluer la faisabilité de l'utilisation des architectures \textbf{GPU} pour la simulation de réseaux de neurones. Il s'agit de déterminer la capacité maximale du système en nombre de \textbf{synapses} prises en charge, tout en précisant les conditions dans lesquelles le calcul sur \textbf{GPU} surpasse une exécution sur \textbf{CPU}. Enfin, les travaux visent à accélérer et à optimiser les performances de simulation en s'appuyant sur nos compétences techniques.

L'approche suivie s'articule en cinq étapes : modélisation \textbf{Hodgkin-Huxley}, couplage par \textbf{synapses}, construction d'un \textbf{connectome}, portage sur \textbf{GPU}, puis \textbf{benchmark} \textbf{CPU}/\textbf{GPU}. L'objectif final est de produire une courbe de \textbf{speedup} en fonction de la taille du réseau, qui montre à partir de quel nombre de neurones le \textbf{GPU} devient indispensable, et dans quelle mesure il l'est.

\section{Le modèle de Hodgkin-Huxley}

Le socle de code fourni repose sur le modèle de \textbf{Hodgkin-Huxley}. Il s'agit d'une modélisation biologique dans laquelle un neurone est assimilé à un circuit électrique, simulant l'activité cérébrale à partir du flux d'ions sodium et potassium. Le modèle décrit l'évolution du \textbf{potentiel membranaire} par un système d'équations différentielles ordinaires (\textbf{EDO}) gouvernant l'ouverture et la fermeture des \textbf{canaux ioniques}. C'est ce formalisme, une \textbf{EDO} par neurone répétée sur l'ensemble du réseau, qui rend le problème naturellement parallélisable.

\section{Problématique}

\textbf{Problématique.} Comment accélérer la simulation de réseaux de neurones de \textbf{Hodgkin-Huxley} couplés lorsque la taille du \textbf{connectome} augmente, et quelle stratégie \textbf{GPU} offre le meilleur gain par rapport à la référence \textbf{CPU}, à résultats équivalents ?

\textbf{La vraie question du projet.} Le laboratoire de M.~Pirog peut-il, grâce aux serveurs Junia équipés de \textbf{RTX~5090}, simuler des réseaux de neurones assez grands pour être utiles à la recherche, dans un temps acceptable ?

\begin{figure}[p]
\centering
\begin{tikzpicture}[
    >={Stealth[length=2.4mm,width=1.9mm]},
    % --- bloc central : la problematique ---
    hub/.style={
        rectangle, rounded corners=5pt, align=center,
        fill=hubbg, draw=hubbg, line width=1pt,
        text=white, inner sep=0pt
    },
    % --- carte d'axe ---
    card/.style={
        rectangle, rounded corners=4pt, align=left,
        line width=1pt, inner sep=8pt, font=\footnotesize,
        text=cardink
    },
    % --- liaisons : un seul gris, la couleur reste aux cartes ---
    link/.style={->, line width=0.9pt, draw=cardline, rounded corners=4pt},
    flow/.style={->, line width=0.9pt, draw=cardline}
]

% =====================================================================
%  Geometrie : 2 colonnes x 3 lignes, tout aligne
% =====================================================================
\def\colL{-3.85}   \def\colR{3.85}
\def\rowA{0}       \def\rowB{-3.6}    \def\rowC{-7.2}
\def\cardw{6.55cm} \def\cardh{2.3cm}

% --- Bloc central ---------------------------------------------------
\node[hub] (HUB) at (0,3.55) {%
  \begin{minipage}{14.8cm}
    \centering\vspace{9pt}
    {\scriptsize\bfseries\color{hubrule}PROBLÉMATIQUE}\\[4pt]
    {\large\bfseries SIMULATION GPU DE RÉSEAUX DE NEURONES}\\[5pt]
    {\color{hubrule}\rule{13.6cm}{0.5pt}}\\[5pt]
    {\small Le GPU (RTX~5090) apporte-t-il un gain suffisant face au CPU ?}
    \vspace{9pt}
  \end{minipage}};

% --- Ligne 1 : POURQUOI  |  COMMENT ----------------------------------
\node[card, draw=whyborder, fill=whybg] (WHY) at (\colL,\rowA) {%
  \begin{minipage}[t][\cardh][t]{\cardw}\cardsetup{whyborder}
    \cardtitle{whyborder}{POURQUOI ?}
    \begin{cbul}
      \item Passer de plusieurs heures (CPU) à quelques minutes.
      \item Simuler de plus grands réseaux de neurones.
      \item Tester plus vite de nouvelles hypothèses.
    \end{cbul}
  \end{minipage}};

\node[card, draw=howborder, fill=howbg] (HOW) at (\colR,\rowA) {%
  \begin{minipage}[t][\cardh][t]{\cardw}\cardsetup{howborder}
    \cardtitle{howborder}{COMMENT ?}
    \begin{cbul}
      \item Code de référence CPU fourni.
      \item Parallélisation sur GPU (RTX~5090).
      \item Benchmark CPU vs GPU.
    \end{cbul}
  \end{minipage}};

% --- Ligne 2 : QUOI  |  ENJEUX ---------------------------------------
\node[card, draw=whatborder, fill=whatbg] (WHAT) at (\colL,\rowB) {%
  \begin{minipage}[t][\cardh][t]{\cardw}\cardsetup{whatborder}
    \cardtitle{whatborder}{QUOI ?}
    \begin{cbul}
      \item Modèle Hodgkin-Huxley (ions Na$^{+}$/K$^{+}$).
      \item Synapses et connectome.
      \item Réseau de neurones couplés.
    \end{cbul}
  \end{minipage}};

\node[card, draw=outborder, fill=outbg] (OUT) at (\colR,\rowB) {%
  \begin{minipage}[t][\cardh][t]{\cardw}\cardsetup{outborder}
    \cardtitle{outborder}{ENJEUX}
    \begin{cbul}
      \item Courbe de speedup selon la taille du réseau.
      \item Seuil à partir duquel le GPU est utile.
      \item Réponse mesurée, pas prédéterminée.
    \end{cbul}
  \end{minipage}};

% --- Ligne 3 : QUI  |  CONTRAINTES -----------------------------------
\node[card, draw=whoborder, fill=whobg] (WHO) at (\colL,\rowC) {%
  \begin{minipage}[t][\cardh][t]{\cardw}\cardsetup{whoborder}
    \cardtitle{whoborder}{QUI ?}
    \begin{cbul}
      \item Laboratoire de M.~Pirog (Junia).
      \item Chercheurs en neurosciences computationnelles.
      \item Étudiants du projet R\&D.
    \end{cbul}
  \end{minipage}};

\node[card, draw=defiborder, fill=defibg] (DEFI) at (\colR,\rowC) {%
  \begin{minipage}[t][\cardh][t]{\cardw}\cardsetup{defiborder}
    \cardtitle{defiborder}{CONTRAINTES}
    \begin{cbul}
      \item Couplage synaptique irrégulier.
      \item Accès mémoire non uniforme.
      \item Part séquentielle inévitable.
    \end{cbul}
  \end{minipage}};

% =====================================================================
%  Liaisons : une fourche depuis le bloc central, puis deux colonnes
% =====================================================================
\draw[link] (HUB.south) -- (0,2.05) -- (\colL,2.05) -- (WHY.north);
\draw[link] (HUB.south) -- (0,2.05) -- (\colR,2.05) -- (HOW.north);

\draw[flow] (WHY.south)  -- (WHAT.north);
\draw[flow] (WHAT.south) -- (WHO.north);
\draw[flow] (HOW.south)  -- (OUT.north);
\draw[flow] (OUT.south)  -- (DEFI.north);

\end{tikzpicture}
\caption{Carte conceptuelle du projet : qui, quoi, comment, pourquoi, enjeux et contraintes.}
\end{figure}

\end{document}
